\section{Mathematical Preliminaries \& System Background}
\label{sec:background}

This section establishes the mathematical foundations, differential geometry, information theory, and hardware principles underpinning the NeuralOps framework.

\subsection{Multi-Head Softmax Attention and Complexity Limits}
Let $X \in \mathbb{R}^{N \times d_{\text{model}}}$ denote an input sequence of $N$ tokens embedded in a $d_{\text{model}}$-dimensional space. Multi-head self-attention projects $X$ into Query ($Q$), Key ($K$), and Value ($V$) matrices across $H$ heads:
\begin{equation}
Q = X W_Q, \quad K = X W_K, \quad V = X W_V,
\end{equation}
where $W_Q, W_K \in \mathbb{R}^{d_{\text{model}} \times d_k}$ and $W_V \in \mathbb{R}^{d_{\text{model}} \times d_v}$. For each head, the scaled dot-product attention computes:
\begin{equation}
A = \operatorname{softmax}\left(\frac{Q K^\top}{\sqrt{d_k}} + M\right) \in \mathbb{R}^{N \times N}, \quad O = A V \in \mathbb{R}^{N \times d_v},
\end{equation}
where $M$ denotes an optional autoregressive causal mask ($M_{ij} = -\infty$ for $j > i$). The explicit materialization of the $N \times N$ affinity matrix $A$ requires $\mathcal{O}(N^2)$ memory and $\mathcal{O}(N^2 d_k)$ floating-point operations (FLOPs), creating a quadratic compute and memory barrier.

\subsection{Linear Attention \& Associative Recurrent State-Space Formulation}
Linear attention replaces the exponential kernel with a decomposable non-negative feature map $\phi: \mathbb{R}^{d_k} \to \mathbb{R}^{m}$ such that $\exp(q_i^\top k_j / \sqrt{d_k}) \approx \phi(q_i)^\top \phi(k_j)$. Exploiting the associativity of matrix multiplication, the causal output at time step $t$ is expressed as:
\begin{equation}
O_t = \frac{\sum_{i=1}^t \left(\phi(Q_t) \phi(K_i)^\top\right) V_i}{\sum_{i=1}^t \left(\phi(Q_t) \phi(K_i)^\top\right) \mathbf{1}} = \frac{\phi(Q_t) \left(\sum_{i=1}^t \phi(K_i)^\top V_i\right)}{\phi(Q_t) \left(\sum_{i=1}^t \phi(K_i)^\top \mathbf{1}\right)} = \frac{\phi(Q_t) S_t}{\phi(Q_t) Z_t + \epsilon},
\end{equation}
where the recurrent associative state matrix $S_t \in \mathbb{R}^{m \times d_v}$ and normalization vector $Z_t \in \mathbb{R}^m$ follow the update rules:
\begin{equation}
S_t = S_{t-1} + \phi(K_t)^\top V_t, \quad Z_t = Z_{t-1} + \phi(K_t)^\top \mathbf{1}.
\end{equation}
This formulation reduces per-token inference time to $\mathcal{O}(m d_v)$ with constant state size $\mathcal{O}(m d_v)$, completely eliminating the growing Key-Value (KV) cache.

\subsection{Truncated Singular Value Decomposition (SVD) \& Eckart-Young Optimality}
For any linear weight matrix $W \in \mathbb{R}^{m \times n}$, the Singular Value Decomposition (SVD) decomposes $W$ as:
\begin{equation}
W = U \Sigma V^\top = \sum_{i=1}^{\min(m,n)} \sigma_i u_i v_i^\top,
\end{equation}
where $U \in \mathbb{R}^{m \times m}$ and $V \in \mathbb{R}^{n \times n}$ are orthogonal matrices ($U^\top U = I_m, V^\top V = I_n$), and $\Sigma = \operatorname{diag}(\sigma_1, \sigma_2, \dots) \in \mathbb{R}^{m \times n}$ contains singular values in non-increasing order ($\sigma_1 \ge \sigma_2 \ge \dots \ge 0$).

\begin{theorem}[Eckart-Young-Mirsky Theorem~\cite{eckart1936approximation}]
Let $W_r = \sum_{i=1}^r \sigma_i u_i v_i^\top = U_r \Sigma_r V_r^\top$ denote the rank-$r$ truncated SVD approximation ($r < \min(m,n)$). For any matrix $B \in \mathbb{R}^{m \times n}$ with $\operatorname{rank}(B) \le r$:
\begin{equation}
\min_{\operatorname{rank}(B) \le r} \|W - B\|_F = \|W - W_r\|_F = \sqrt{\sum_{i=r+1}^{\min(m,n)} \sigma_i^2},
\end{equation}
and $\min_{\operatorname{rank}(B) \le r} \|W - B\|_2 = \|W - W_r\|_2 = \sigma_{r+1}$.
\end{theorem}
The Eckart-Young theorem guarantees that truncated SVD provides the strictly optimal low-rank projection in both Frobenius and spectral operator norms.

\subsection{Riemannian Geometry on the Stiefel Manifold $\operatorname{St}(p, n)$}
The Stiefel manifold $\operatorname{St}(p, n)$ represents the set of all orthonormal $p$-frames in $\mathbb{R}^n$:
\begin{equation}
\operatorname{St}(p, n) = \left\{ W \in \mathbb{R}^{n \times p} \mid W^\top W = I_p \right\}, \quad (p \le n).
\end{equation}
Endowed with the canonical Euclidean metric $\langle \xi, \eta \rangle = \operatorname{tr}(\xi^\top \eta)$, the tangent space $\mathcal{T}_W \operatorname{St}(p, n)$ at point $W$ is defined by:
\begin{equation}
\mathcal{T}_W \operatorname{St}(p, n) = \left\{ \xi \in \mathbb{R}^{n \times p} \mid W^\top \xi + \xi^\top W = 0 \right\}.
\end{equation}
Given an unconstrained Euclidean gradient $G = \nabla_W \mathcal{L} \in \mathbb{R}^{n \times p}$, its orthogonal projection onto the tangent space $\mathcal{T}_W \operatorname{St}(p, n)$ is given by the Riemannian gradient:
\begin{equation}
\operatorname{grad}_R \mathcal{L} = \Pi_{\mathcal{T}_W}(G) = G - W \operatorname{sym}(W^\top G) = G - \frac{1}{2} W \left(W^\top G + G^\top W\right).
\end{equation}
To map a tangent vector $\xi \in \mathcal{T}_W \operatorname{St}(p, n)$ back onto the manifold, a retraction $\mathcal{R}_W: \mathcal{T}_W \operatorname{St}(p, n) \to \operatorname{St}(p, n)$ is required. NeuralOps implements the efficient QR retraction:
\begin{equation}
\mathcal{R}_W(\xi) = \operatorname{qf}(W + \xi),
\end{equation}
where $\operatorname{qf}(A)$ denotes the $Q$ factor of the QR decomposition $A = Q R$ with positive diagonal elements on $R$.

\subsection{Hardware Memory Hierarchy \& The Roofline Performance Model}
Modern GPUs organize compute and memory in a strict hierarchical architecture:
\begin{itemize}
    \item \textbf{Global High-Bandwidth Memory (HBM):} High capacity ($24$--$80\text{ GB}$) but high access latency ($\sim 200$--$400\text{ cycles}$) and bounded bandwidth ($1$--$3.35\text{ TB/s}$).
    \item \textbf{On-Chip SRAM \& Register Files:} Extremely low latency ($\sim 1$--$2\text{ cycles}$) and ultra-high aggregate bandwidth ($>30\text{ TB/s}$), but limited capacity ($100$--$228\text{ KB}$ per Streaming Multiprocessor).
\end{itemize}

The Roofline performance model~\cite{williams2009roofline} characterizes achievable throughput $P$ (TFLOPS) as a function of operational arithmetic intensity $\mathcal{I}$ ($\text{FLOPs}/\text{Byte}$ transferred from HBM):
\begin{equation}
P(\mathcal{I}) = \min\left(\pi_{\text{peak}}, \mathcal{I} \cdot \beta_{\text{peak}}\right),
\end{equation}
where $\pi_{\text{peak}}$ is peak Tensor Core compute throughput and $\beta_{\text{peak}}$ is peak memory bandwidth. When $\mathcal{I} < \mathcal{I}^* = \pi_{\text{peak}} / \beta_{\text{peak}}$, the operation is memory-bandwidth bound. Bypassing intermediate HBM roundtrips by fusing operations directly within SRAM increases $\mathcal{I}$, pushing execution toward the compute-bound ceiling $\pi_{\text{peak}}$.

\subsection{Information-Theoretic Spectral Entropy}
Given the singular value spectrum $\Sigma = (\sigma_1, \dots, \sigma_r)$ of a projection weight matrix, we define the normalized singular probability distribution $p_i = \sigma_i / \sum_{j=1}^r \sigma_j$. The spectral entropy $H(\Sigma)$ is defined as:
\begin{equation}
H(\Sigma) = -\sum_{i=1}^r p_i \ln p_i.
\end{equation}
A high spectral entropy ($H(\Sigma) \approx \ln r$) indicates an isotropic, uniform distribution of information across all singular subspaces, whereas low entropy ($H(\Sigma) \ll \ln r$) reflects strong low-rank concentration along a few dominant principal directions.
